\documentclass[11pt,a4paper]{article}

%% PACHETE MATEMATICE
\usepackage{amsmath,amssymb,amsfonts}
\usepackage{amsthm}
\usepackage{mathpazo}
\usepackage{eulervm}
\usepackage{enumerate}
\usepackage{color}
\usepackage{graphicx}
\usepackage[all]{xy}
\usepackage{xcolor}
\usepackage{framed}
\usepackage[unicode]{hyperref}
\setcounter{MaxMatrixCols}{20}
\usepackage{multicol}
% \usepackage[romanian]{babel}


%% ASPECT
\linespread{1.05}
\usepackage[T1]{fontenc}
\usepackage[utf8x]{inputenc}
\usepackage[margin=2cm]{geometry}
% \usepackage[color]{showkeys}
\usepackage{anyfontsize}
\usepackage{titlesec}
\titleformat*{\section}{\large\bfseries}

\usepackage{fancyhdr}
\pagestyle{fancy}
\rhead{\small \color{gray} Notițe de Adrian Manea}
\lhead{\small \color{gray} IntProb-FIA, 2017---2018, L. Lemnete \& M. Olteanu}


\usepackage[autostyle = false, style = german]{csquotes}
\usepackage[romanian]{babel}
\newcommand\qq{\enquote}

%% COMENZILE MELE
\renewcommand*{\proofname}{Dem.:}
\newcommand\bb{\mathbb}
\newcommand\dr{\mathrm}
\newcommand\kal{\mathcal}
\newcommand\fk{\mathfrak}
\newcommand\op{\oplus}
\newcommand\ot{\otimes}
\newcommand\ds{\displaystyle}
\newcommand\id{\indent}
\newcommand\nid{\noindent}
\newcommand\seq{\subseteq}
\newcommand\sneq{\subsetneq}
\newcommand\speq{\supseteq}
\newcommand\spneq{\supsetneq}
\newcommand\rar{\rightarrow}
\newcommand\Rar{\Rightarrow}
\newcommand\xrar{\xrightarrow}
\newcommand\lar{\leftarrow}
\newcommand\Lar{\Leftarrow}
\newcommand\xlar{\xleftarrow}
\newcommand\lrar{\leftrightarrow}
\newcommand\Lrar{\Leftrightarrow}
\newcommand\ideal{\trianglelefteq}
\newcommand\ai{\text{ a.\^{i}. }}
\newcommand\sk{(\textit{Schi\c{t}\u{a}}) }
\newcommand\rhup{\rightharpoonup}
\newcommand\rhdn{\rightharpoondown}
\newcommand\lhup{\leftharpoonup}
\newcommand\lhdn{\leftharpoondown}
\newcommand\vid{\emptyset}
\newcommand\wh{\widehat}
\newcommand\ol{\overline}
\newcommand\NN{\mathbb{N}}
\newcommand\RR{\mathbb{R}}
\newcommand\ZZ{\mathbb{Z}}
\newcommand\QQ{\mathbb{Q}}
\newcommand\CC{\mathbb{C}}

%% STILURILE MELE
\newtheoremstyle{dotless}{}{}{\itshape}{}{\bfseries}{:}{ }{}

\theoremstyle{dotless}
\newtheorem{theorem}{Teorem\u{a}}[section]
\newtheorem{proposition}{Propozi\c{t}ie}[section]
\newtheorem{lemma}{Lem\u{a}}[section]
\newtheorem{corollary}{Corolar}[section]
\newtheorem{exercise}{Exerci\c{t}iu}

\newtheoremstyle{dotlessdr}{}{}{}{}{\bfseries}{:}{ }{}
\theoremstyle{dotlessdr}
\newtheorem{example}{Exemplu}[section]
\newtheorem{definition}{Defini\c{t}ie}[section]
\newtheorem{remark}{Observa\c{t}ie}[section]




\begin{document}

\begin{center}
  {\large\textbf{Seminar 8 \\ Recapitulare parțial}}
\end{center}

\vspace{1cm}

\textbf{Șirul aproximațiilor succesive}

1. Să se aproximeze cu o eroare mai mică decît $ 10^{-3} $ soluția reală
a ecuației $ x^3 + 4x - 1 = 0 $.

\emph{Indicație:} $ f(x) = \frac{1}{x^2 + 4} $ sau $ g(x) = \frac{1}{4}(1 - x^3) $.

\vspace{1cm}

2. Aceeași cerință pentru $ x^3 + 12x - 1 = 0 $.

\emph{Indicație:} $ f(x) = \frac{1}{x^2 + 12} $.

\vspace{1cm}

3. Aceeași cerință pentru: $ x^5 + x - 15 = 0 $.

\emph{Indicație:} $ f(x) = \sqrt[5]{15 - x} $.

\vspace{2cm}

%%%%%%%%%%%%%%%%%%%%%%%%%%%%%%%%%%%%%%%%%%%%%%%%%%%%%%%%%%%%%%%%%%%%%%

\textbf{Integrale improprii}

4. Să se calculeze, folosind funcțiile $\Gamma$ și $B$, integralele:
\begin{enumerate}[(a)]
\item $ \ds\int_0^\infty e^{-x^p} dx, p > 0 $;
\item $ \ds\int_0^{\infty} \frac{x^{\frac{1}{4}}}{(x + 1)^2} dx $;
\item $ \ds\int_0^\infty \frac{dx}{x^3 + 1} dx $;
\item $ \ds\int_0^{\frac{\pi}{2}} \sin^p x \cos^q x dx, p > -1, q > - 1 $;
\item $ \ds\int_0^1 x^{p+1}(1-x^m)^{q - 1} dx, p, q, m > 0 $;
\item $ \ds\int_0^\infty x^p e^{-x^q} dx, p > -1, q > 0 $;
\item $ \ds\int_0^1 \ln^p x^{-1} dx, p > -1 $;
\item $ \ds\int_0^1 \frac{dx}{(1 - x^n)^{\frac{1}{n}}}, n \in \NN $;
\item $ \ds\int_1^e \frac{1}{x} \ln^3 x \cdot (1 - \ln x)^4 dx $;
\item $ \ds\int_0^1 \sqrt{\ln \frac{1}{x}} dx $;
\item $ \ds\int_0^\infty \frac{x^2}{(1 + x^4)^2} dx $;
\item $ \ds\int_{-1}^\infty e^{-x^2 - 2x + 3} dx $;
\end{enumerate}

\vspace{2cm}

%%%%%%%%%%%%%%%%%%%%%%%%%%%%%%%%%%%%%%%%%%%%%%%%%%%%%%%%%%%%%%%%%%%%%% 


\textbf{Integrale curbilinii}

5. Să se calculeze următoarele integrale curbilinii de speța întîi:
\begin{enumerate}[(a)]
\item $ \int_C x ds$, unde $C: y = x^2, x \in [0, 2] $;
\item $ \int_C y^5 ds$, unde $C: x = \frac{y^4}{4}, y \in [0, 2] $;
\item $ \int_C x^2 ds$, unde $C: x^2 + y^2 = 2, x, y \geq 0 $;
\item $ \int_C yds$, unde $C: x(t) = \ln(\sin t) - \sin^2 t $,
  $y(t) = \frac{1}{2} \sin 2t, t \in \big[ \frac{\pi}{3}, \frac{\pi}{2} \big]$;
\item $ \int_C xyds$, unde $C: x(t) = |t|, y(t) = \sqrt{1 - t^2}$, $t \in [-1, 1] $;
\item $ \int_C |x - y| ds$, unde $C: x(t) = |\cos t|, y(t) = \sin t$, $t \in [0, \pi] $.
\end{enumerate}

\vspace{1cm}

6. Să se calculeze următoarele integrale curbilinii de speța a doua:
\begin{enumerate}[(a)]
\item $ \int_C \frac{dx}{2a + y} - \frac{dy}{a + x}$, unde $C$ este o curbă
  simplă formată din porțiunea din cercul $x^2 + y^2 + 2ay = 0$ ($a > 0$),
  pentru care $x + y \geq 0$, iar capătul de pornire este $A(a, -a)$;
\item $\int_C xdy - ydx$, unde $C: x^2 + y^2 = 4, y = x \sqrt{3}, x \geq 0, %
  y = \frac{x}{\sqrt{3}}$;
\item $\int_\gamma (x + y) dx + (x - y) dy$, unde domeniul este:
  \[
    \gamma = \{(x, y) \mid x^2 + y^2 = 4, y \geq 0\}.
  \]
\item $\int_\gamma \frac{y}{x + 1} dx + dy$, unde $\gamma$ este triunghiul cu
  vîrfurile $A(2, 0), B(0, 0), C(0, 2)$;
\item $\int_\gamma xdy - ydx$, unde:
  \[
    \gamma = \{(x, y) \mid \frac{x^2}{a^2} + \frac{y^2}{b^2} = 1 \}.
  \]
\end{enumerate}

\vspace{1cm}

7. Să se calculeze $\int_\gamma ydx + xdy$, pe un drum de la $A(2, 1)$
la $B(1, 3)$.

\vspace{2cm}

%%%%%%%%%%%%%%%%%%%%%%%%%%%%%%%%%%%%%%%%%%%%%%%%%%%%%%%%%%%%%%%%%%%%%%


\textbf{Integrale duble și triple}

8. Calculați volumele mulțimilor $ \Omega $, mărginite de suprafețele de ecuații:
\begin{enumerate}[(a)]
\item $ 2x^2 + y^2 + z^2 = 1, \quad 2x^2 + y^2 - z^2 = 0, z \geq 0 $;
\item $ z = x^2 + y^2 - 1, \quad z = 2 - x^2 - y^2 $;
\item $ z = 4 - x^2 - y^2, \quad 2z = 5 + x^2 + y^2 $;
\item $ x^2 + y^2 = 1, \quad z^2 = x^2 + y^2, \quad z \geq 0 $.
\end{enumerate}

\vspace{1cm}

9. Să se calculeze, trecînd la coordonate polare, integralele:
\begin{enumerate}[(a)]
\item $ \ds\iint_D e^{x^2 + y^2} dxdy$, unde $D = \{(x, y) \in \RR^2 \mid x^2 + y^2 \leq 1 \} $;
\item $ \ds\iint_D 1 + \sqrt{x^2 + y^2} dxdy$, unde $D = \{(x, y) \in \RR^2 \mid %
  x^2 + y^2 - y \leq 0, x \geq 0\} $;
\item $ \ds\iint_D \ln(1 + x^2 + y^2) dxdy $, unde $D$ este mărginit de curbele de ecuații:
  \[
    \begin{cases}
      x^2 + y^2 &= e^2 \\
      y &= x \sqrt{3} \\
      x &= y\sqrt{3} \\
      x &\geq 0
    \end{cases}
  \]
\end{enumerate}

\vspace{1cm}

10. Calculați integralele duble, găsind explicit domeniul de definiție (prin intersecția
curbelor care îl definesc):
\begin{enumerate}[(a)]
\item $ \ds\iint_D xydxdy$, cu $D: xy = 1, \quad x + y = \frac{5}{2} $;
\item $ \ds\iint_D \frac{1}{\sqrt{x}} dxdy$, cu $D: y^2 \leq 8x, y \leq 2x, y + 4x \leq 24 $;
\item $ \ds\iint_D (1 - y) dxdy$, cu $D: x^2 + (y-1)^2 \leq 1, y \leq x^2, x \geq 0 $.
\end{enumerate}


\end{document}